% !TeX root = RJwrapper.tex
\title{Working with table environments in texor}
\author{by Abhishek Ulayil}

\maketitle

\abstract{
This is a small sample article to demonstrate usage of \CRANpkg{texor} to convert table environments.
}

\section{Introduction}
Tables are commonly used in R Journal articles to display data in a tabular format. However, there are differences in the way tables are handled by LaTeX and HTML.
LaTeX tables have more customization and are usually optimized for printing, whereas the web articles need tables optimized for varying sizes of media.
Pandoc converts most of the tables somewhat easily, but is unable to do well with table customization packages and complex tables.
The \CRANpkg{texor} packages uses the pandoc extensions \verb|simple_tables| and \verb|pipe_tables| to tackle these difficult cases.

\section{Generic tables}
Normal LaTeX tables will be converted just fine, even if they include a some math and other elements.
However, note that the conversion will be done to traditional or pandoc-style markdown rather than R Markdown with the \verb|knitr::kable(..)| function. Also any custom LaTeX commands or specific font characters will
not be supported.
\begin{table}[htbp]
\centering
\begin{tabular}{l | llll }
 \hline
 Graphics Format & LaTeX & Markdown & R Markdown & HTML \\
 \hline
 PNG       & Yes & Yes & Yes & Yes \\
 JPG       & Yes & Yes & Yes & Yes \\
 PDF       & Yes & No & No & No \\
 SVG       & No & Yes & Yes & Yes \\
 Tikz      & Yes & No & Yes & No \\
 Algorithm & Yes & No & No & No \\
\hline
\end{tabular}
\caption{Image format support in various markup/typesetting languages.}
\label{table:1}
\end{table}

\section{Multicolumn tables}
\begin{table*}[htbp]
\begin{center}
\begin{tabular}{l | llll }
 \hline
 \multicolumn{1}{c |}{EXAMPLE} & \multicolumn{2}{c}{$X$} &
\multicolumn{2}{c}{$Y$} \\
 \hline
  & 1 & 2 & 1 & 2 \\
 EX1  & X11 & X12 &  Y11  & Y12 \\
 EX2  & X21 & X22 &  Y21  & Y22 \\
 EX3  & X31 & X32 &  Y31  & Y32 \\
 EX4  & X41 & X42 &  Y41  & Y42\\
 EX5  & X51 & X52 &  Y51  & Y52 \\
\hline
\end{tabular}
\caption{An example table using \code{multicolumn}}
\label{table:2}
\end{center}
\end{table*}

Table~\ref{table:2} illustrates a table that uses the \code{multicolumn} command, which the \CRANpkg{texor} package can handle through subroutines and pre-processing steps to transform the LaTeX code. 

\begin{itemize}
    \item The Stream Editor modifies the \verb|\multicolumn{..}| to \verb|\multicolumnx{..}|
    \item Then a LaTeX macro redefines the \verb|\multicolumnx{..}| to \verb|\multicolumn{---}|, which is accepted by pandoc.
    \item Finally pandoc reads the table and transforms it to markdown.
\end{itemize}

Also note that the stream editor is used to rename \verb|table*| environment to \verb|table| environment because the HTML format is single column, so the asterisk indicating that the table should be drawn over the full width of the page is redundant in this case.

\section{Complex tables}
A complex table with various other elements like figures, math, code and so on, are also supported
by \CRANpkg{texor}.  
\begin{table}[htbp]
\centering
\begin{tabular}{|l | lll |}
 \hline
 Inline Format & LaTeX Support & Web Support &  Rendering \\
 \hline
 Text       & Yes & Yes  & \emph{Hello} \\ \hline
 Image       & Yes & Yes & \begin{minipage}{0.25\textwidth}
    \centering
    \includegraphics[width=1\textwidth]{Rlogo-5.png}
    \end{minipage} \\ \hline
 CodeBlock      & Yes & Yes & \begin{minipage}{0.45\textwidth}
\vspace{1mm}
\begin{example}
x <- 1:100
y <- dbinom(x,100,prob = 0.5)
plot(x,y)
\end{example}
    \end{minipage} \\ \hline
 Math      & Yes & Yes & $e = m c^2 $ \\ \hline
 Link      & Yes & Yes &  \href{www.google.com}{Google} \\ \hline
 Nested Table & Yes & No & NaN \\
\hline
\end{tabular}
\caption{Image format support in various markup/typesetting languages}
\label{table:3}
\end{table}

\section{Wide tables}

\begin{widetable}
\centering
\begin{tabular}{lllllllll}
\toprule
index  &  A  &  B     & C   &  D  &  E  &  F  &  G  & H  \\
\midrule
1 & 359.00 & NaN & 5796.00 & 0.00 & 16.14 & 1.00 & NaN & 0.00\\
2 & 25.73 & 0.00 & 1029.20 & NaN & 40.00 & 0.68 & 0.00 & 0.00\\
2.1 & 26.26 & 0.00 & 13.40 & 0.00 & 2.14 & 0.68 & 0.00 & NaN\\
2.2 & 32.06 & 20.06 & 47.64 & 0.04 & 1.80 & 0.68 & 0.01 & NaN\\
2.3 & 51.94 & 420.27 & 21.17 & 0.20 & 1.77 & 0.74 & 0.05 & NaN\\
2.4 & 40.62 & 30.44 & \phantom{000}0.90 & 0.57 & 1.44 & 1.31 & 0.24 & NaN\\
 & \ldots & \ldots&  \ldots& \ldots &  \ldots& \ldots& \ldots & \ldots \\
\bottomrule
\end{tabular}
\begin{tabular}{llllllll}
\toprule
index &    I  &  J  &  K  &  L  &  M  &  N & O \\
\midrule
1 & 0.00 & $-$1.73 & 0.00 & 0.00 & 0.00 & 0.00 & NaN\\
2 & 0.21 & $-$33.41 & 0.00 & 0.11 & 0.00 & NaN & $-$197.85\\
2.1 & 0.24 & $-$24.00 & 0.00 & 0.15 & NaN & 0.06 & $-$70.46\\
2.2 & 0.6 & $-$19.42 & 0.00 & 0.15 & 0.00 & 0.11 & $-$16.48\\
2.3 & 0.75 & $-$ 31.77 & 0.00 & 0.18 & 0.01 & 0.37 & $-$0.82\\
2.4 & 0.26 & $-$1.89 & 0.1 & 0.55 & 0.70 & 0.22 & $-$6.55 \\
  &\ldots &  \ldots&  \ldots& \ldots & \ldots & \ldots & \ldots \\
\bottomrule
\end{tabular}
\caption{\code{A dummy research data}}
\label{table:4}

\end{widetable}

Widetable environments are treated specially when converted by the \CRANpkg{texor} package. 
They get their own block and separate numbering (although references are shared by figures).

\begin{widetable}
\centering %
\begin{tabular}{llrr}
\toprule 
Package  & Commits  & Version  & Last Updated \tabularnewline
\toprule 
\CRANpkg{texor} & 260 & 1.1.0  & 28-Jul-2023 \tabularnewline
\CRANpkg{rebib} & 76 & 0.2.4  & 29-Jul-2023 \tabularnewline
\CRANpkg{rjtools}  & 314 & 1.0.11  & 30-Jul-2023 \tabularnewline
\CRANpkg{rmarkdown} & 3189 & 2.23  & 31-Jul-2023 \tabularnewline
\toprule 
\end{tabular}
\caption{A dummy summary of a few CRAN packages}
\label{table:5}

\end{widetable}


Here is a reference to Table~\ref{table:4} and Table~\ref{table:5}.
\section{Summary}

In summary the \CRANpkg{texor} package supports:
\begin{itemize}
\item Some common table environments.
\item Wide tables.
\item Some tables with \code{multicolumn} commands.
\item Environments such as \code{figure}/\code{code} in tables.
\end{itemize}



\begin{thebibliography}{2}
    \providecommand{\natexlab}[1]{#1}
    \providecommand{\url}[1]{\texttt{#1}}
    \expandafter\ifx\csname urlstyle\endcsname\relax
      \providecommand{\doi}[1]{doi: #1}\else
      \providecommand{\doi}{doi: \begingroup \urlstyle{rm}\Url}\fi

\bibitem[Krewinkel, Lucero (2023)]{pandoc}
A.~ Krewinkel and A.~ Lucero
\newblock pandoc 3.0 Release notes
\newblock \emph{pandoc} \penalty0 2023
\newblock URL \url{https://pandoc.org/releases.html}

\end{thebibliography}


\address{%
Abhishek Ulayil\\
Student, Institute of Actuaries of India\\%
Mumbai, India\\
ORCiD: 0009-0000-6935-8690\\
}
